% ======================= COMPILATION NOTICE =======================
% Before compiling, place ge2340_project.sty and all three template
% files (phase1_template.tex, phase2_template.tex, and
% contribution_template.tex) in the same directory.
% Keep the style file named ge2340_project.sty so that the templates
% can load it correctly.
% ==================================================================

\documentclass{article}

% ======================= REQUIRED STYLE =======================
\PassOptionsToPackage{numbers}{natbib}
\usepackage[preprint]{ge2340_project}

\usepackage[utf8]{inputenc}
\usepackage[T1]{fontenc}
\usepackage{hyperref}
\usepackage{url}
\usepackage{booktabs}
\usepackage{amsmath,amssymb}
\usepackage{graphicx}
\usepackage{microtype}
\usepackage{xcolor}
% \usepackage{tabularx}

% ======================== ASSIGNMENT CONSTRAINTS ========================
% Phase 1 is GROUP WORK. Each group has 3--4 students.
% The proposal must be in English and submitted as a PDF.
% The complete PDF must not exceed 2 pages, including figures and references.
% Do NOT rename, remove, or reorder the required section titles below.
% Replace only the bracketed placeholders and remove instructional text before submission.
% =======================================================================

\title{Phase 1: AI Project Proposal\\
\large Email Content Summarization \& Scheduler}

% Use the fourth author block only for a four-member group.
\author{
  FENG Yuxiang (Student ID: 58520850)\\
  \texttt{yuxiafeng2-c@my.cityu.edu.hk}
  \And 
  SOETANTO Darrel Edbert (Student ID: 58586617)\\
  \texttt{dsoetanto2-c@my.cityu.edu.hk}
  \AND
  CHEN Lehan (Student ID: 58520456)\\
  \texttt{lehanchen2-c@my.cityu.edu.hk}
  \And
  TANG Shun Chong Adrian (Student ID: 60313420)\\
  \texttt{adriatang6-c@my.cityu.edu.hk}
}

\begin{document}
\maketitle

\noindent\textbf{Project topic:} Students struggle to schedule their timetables well through AIMS.\\
\textbf{Target users/stakeholders:} Students using AIMS, particularly students with large flexibility in course scheduling.

% ======================= Required sections (do NOT rename) =======================

\section{Problem Definition \& Scope}
\subsection{Problem}

The university students will receive a large volume of emails everyday, which contains assignment / quiz / exam announcement, event invitation, security alert, bank statement notification, promotion, etc. Spending time on checking these emails can be annoying. And the performance of filter function by email apps like Gmail and Outlook is not good enough. After checking these emails, make it to the schedule can also be another issues to the university students.

\subsection{Scope}

The project will contains classification of the emails, summarize an email from paragraph form to bullet-point form, integrate different emails to generate a schedule. But we will make an agent to help users to write emails.

\section{AI Task Formulation}

Task: Rank a set of timetables based on selected selection criteria, including, but not necessarily limited to, time of day, day of week, and consecutiveness

Inputs: List of occupied schedule slots / Slot in schedule's status as occupied / unoccupied (6 days of the week × 13 possible hours per day = 78 binary inputs)

Outputs: Scores based on various criteria (eg: “earliness”, “lateness”, “consecutiveness”, etc; one output per criteria, potentially with an additional final layer for an overall aggregate score)

Constraints: Scores outputted from 0-100. Model will only be fed valid timetables (ie. without collisions)

Evaluation: Judge accuracy of suggestions; Scores should agree with the opinion of users.

Data source: Survey students for their opinion on how a set of generated timetables should be scored. No privacy, consent, or access considerations needed.

\section{Proposed AI Approaches}

\subsection{AI Agent}
  build an agent with several tools, including a email reader, summarizer, report generator and scheduler.  And the task can be timely triggered by the agent based on the email content.

\subsection{Reinforcement Learning}
  The agent will treat the received emails as a environment and learn to take actions based on the user feedback.

\subsection{Naive Bayes}
  Use a naive bayes classifier  based on the words frequence and the context of the email. 
\subsection{Decision Tree}
  Use a decision tree classifier based on special words of the email. 
\subsection{Hard-code Rule-based pipeline}
  Juest detecting the class of the rules by predefined grammar and words rules.


\section{Comparative Analysis}
  AI agent method is the best as it naturally supports the multi-step workflow of email processing.




\section{Expected Impact \& Project Feasibility}
If the project is successful, it will help students to manage their time more efficiently and reduce the time spent on checking emails. It can also detecting the fishing and spam emails to avoid potential harm.This project combine the email summarization and the schedule generation, which is very novel and useful.The data need is the email content and the user feedback, which is easy to collect. We plan to let user install the agent application and collect data after the user permission. Other pre-trained model will be trained based on the open-source dataset. The project is feasible and can be completed within the time frame. The biggest risk is the user privacy, but we will not collect any personal information if the user provides consent.


% ======================= References (included in page limit) =======================
% Cite sources in the text and replace the sample entry below.
\begin{thebibliography}{9}
\small  
\bibitem{sample}
Lance Le, ``A Sample Title,'' \emph{A sample journal}, 2023.

\bibitem{citekey} 
Author Maple, ``Rubik's Cube Methods,'' \emph{Journal of Cube}, 2021.
\end{thebibliography}

% Course policy reminder: LLMs may be used to refine and polish writing,
% but may not be used to generate the submission text from scratch.

\end{document}
